%% ============================================================
%% preamble/boxes.tex — semantic callout family (module: boxes)
%%
%% Purpose      : the five authoring-contract callouts —
%%                keyidea / example / warning / sidebar / definition —
%%                as tcolorbox environments. Each takes ONE optional
%%                keyval argument, e.g.:
%%                  \begin{warning}[title={Check the gutter},importance=high]
%%                Keys: any tcolorbox key, plus the template's
%%                importance=low|medium|high (frame emphasis scaling).
%% Dependencies : tcolorbox (packages.tex) + skins/breakable libraries
%%                (loaded here), colors.tex (callout-* semantics),
%%                fonts.tex (sans for titles). Loaded only when
%%                \ifBookModuleBoxes is true (book.yaml modules.boxes).
%% Provenance   : complexity-book's callout family, incl. the smart
%%                page-breaking discipline (`lines before break=3`
%%                keeps a box from leaving a 1-2 line stub before a
%%                page break) — docs/research/complexity-book.md.
%% ============================================================

\tcbuselibrary{skins,breakable}

% Icons for the math/logic callouts (module: boxes-math). fontawesome5
% ships with TeX Live; icons differentiate the four boxes in BOTH the
% color edition and grayscale POD (where the hues collapse). The EPUB
% converter reproduces the icon from a class-based CSS ::before, so the
% glyph survives conversion independently of the LaTeX title text.
\usepackage{fontawesome5}
\newcommand{\iconProofKit}{\faDraftingCompass}
\newcommand{\iconTryIt}{\faPencilRuler}
\newcommand{\iconGoingDeeper}{\faLayerGroup}
\newcommand{\iconArchive}{\faScroll}

% ----------------------------------------------------------------------
% Shared skeleton. Design language: quiet full hairline + a heavier
% LEFT rule carrying the semantic color; title as a small bold sans
% line inside the tint, not a filled banner (banner fills go muddy on
% POD grayscale). `enhanced jigsaw` is the skin that breaks cleanly.
% ----------------------------------------------------------------------
\tcbset{
  bookcallout/.style={
    enhanced jigsaw,
    breakable,
    lines before break=3,     % never orphan a 1-2 line box stub at a page break
    boxrule=0.4pt,
    leftrule=2.5pt,           % semantic color lives here (importance rescales it)
    arc=1pt,                  % near-square: print-friendly, no halftone halos
    left=8pt, right=8pt, top=6pt, bottom=6pt,
    before skip=0.9\baselineskip,
    after skip=0.9\baselineskip,
    fonttitle=\normalfont\sffamily\bfseries\small,
    toptitle=5pt, bottomtitle=2pt,
    titlerule=0pt,            % title sits in the tint, no separator bar
  },
  % importance= — how loudly the frame speaks. Scales rule weights
  % (and, at high, the title weight) but never changes the hue, so a
  % chapter full of callouts still reads as one family.
  importance/.is choice,
  importance/low/.style={leftrule=1pt, boxrule=0.3pt},
  importance/medium/.style={leftrule=2.5pt, boxrule=0.4pt},
  importance/high/.style={leftrule=5pt, boxrule=0.8pt,
    fonttitle=\normalfont\sffamily\bfseries},
}

% Each environment: default title = its semantic name (override with
% title={...}); colors bound to the callout-* semantic tokens so
% \GrayscaleMode retheming happens upstream in colors.tex.
\newtcolorbox{keyidea}[1][]{bookcallout,
  title={Key Idea},
  colframe=callout-keyidea-frame, colback=callout-keyidea-bg,
  coltitle=callout-keyidea-frame, colbacktitle=callout-keyidea-bg,
  #1}

\newtcolorbox{example}[1][]{bookcallout,
  title={Example},
  colframe=callout-example-frame, colback=callout-example-bg,
  coltitle=callout-example-frame, colbacktitle=callout-example-bg,
  #1}

\newtcolorbox{warning}[1][]{bookcallout,
  title={Warning},
  colframe=callout-warning-frame, colback=callout-warning-bg,
  coltitle=callout-warning-frame, colbacktitle=callout-warning-bg,
  #1}

\newtcolorbox{sidebar}[1][]{bookcallout,
  title={Aside},
  colframe=callout-sidebar-frame, colback=callout-sidebar-bg,
  coltitle=callout-sidebar-frame, colbacktitle=callout-sidebar-bg,
  #1}

\newtcolorbox{definition}[1][]{bookcallout,
  title={Definition},
  colframe=callout-definition-frame, colback=callout-definition-bg,
  coltitle=callout-definition-frame, colbacktitle=callout-definition-bg,
  #1}

% ----------------------------------------------------------------------
% Math/logic callouts (module: boxes-math). Distinct hues (colors.tex)
% + a leading icon so a Proof Kit, a Try It, a Going Deeper aside, and a
% From-the-Archive document are instantly told apart — in colour, and in
% grayscale where the hues collapse but the icons remain. Default titles
% carry the icon; to use a custom title, prepend the matching \icon...
% macro, e.g. title={\iconProofKit\ Proof Kit: Cantor's Theorem}.
% ----------------------------------------------------------------------
\newtcolorbox{proofkit}[1][]{bookcallout,
  title={\iconProofKit\hspace{0.5em}Proof Kit},
  colframe=callout-proofkit-frame, colback=callout-proofkit-bg,
  coltitle=callout-proofkit-frame, colbacktitle=callout-proofkit-bg,
  #1}

\newtcolorbox{tryit}[1][]{bookcallout,
  title={\iconTryIt\hspace{0.5em}Try It},
  colframe=callout-tryit-frame, colback=callout-tryit-bg,
  coltitle=callout-tryit-frame, colbacktitle=callout-tryit-bg,
  #1}

\newtcolorbox{goingdeeper}[1][]{bookcallout,
  title={\iconGoingDeeper\hspace{0.5em}Going Deeper},
  colframe=callout-goingdeeper-frame, colback=callout-goingdeeper-bg,
  coltitle=callout-goingdeeper-frame, colbacktitle=callout-goingdeeper-bg,
  #1}

\newtcolorbox{archive}[1][]{bookcallout,
  title={\iconArchive\hspace{0.5em}From the Archive},
  colframe=callout-archive-frame, colback=callout-archive-bg,
  coltitle=callout-archive-frame, colbacktitle=callout-archive-bg,
  #1}
